\section{Reusable statistics and proof relations}\label{sec:design}
\subsection{Audit task and classical fixed-center identity}\label{sec:background}
Let $X=(x_1,\ldots,x_N)$ be an ordered private catalog and $\mathcal C=\{c_0,\ldots,c_{k-1}\}$ public centers. For a public sorted, distinct index set $I$, the audit computes private first-index nearest-center assignments and the public sum
\begin{equation}
\Phi_I(\mathcal C)=\sum_{i\in I}\min_j\|x_i-c_j\|^2.
\end{equation}
The original row commitment, the approved representation and all verification keys must be authenticated to the verifier. Proving that the centers were obtained by $k$-means training is a separate task.

Choose anchor $c_0$ and $U\in\mathbb R^{d\times a}$ with orthonormal columns whose span contains the deployed center differences. Write
\begin{align}
w(x)&=x-c_0,& h(x)&=U^\top w(x),\nonumber\\
r(x)&=(I-UU^\top)w(x),& e(x)&=\|r(x)\|^2.
\end{align}
We state the classical identity to define the relation that will be implemented \cite{feldman2013tiny,feldman2020tiny,cohen2015dimensionality}.
\begin{proposition}[Fixed-center decomposition]\label{prop:head-energy}
For every deployed center $c_j$,
\begin{equation}\label{eq:head-energy}
\|x-c_j\|^2=\|h(x)-h(c_j)\|^2+e(x).
\end{equation}
Thus $h(x)$ determines the nearest-center index, and $(h(x),e(x))$ determines every fixed-center distance.
\end{proposition}
\begin{proof}
Since $r(c_j)=0$, the decomposition $x-c_j=U(h(x)-h(c_j))+r(x)$ is orthogonal. Apply Pythagoras. The same energy is added for every candidate, including ties.
\end{proof}
The deployed span has rank at most $k-1$, so at most $k$ scalar values per record suffice in real arithmetic. Energy is a squared norm with a different range from a head coordinate; scalar counts are not byte counts.

\subsection{Integer derivation and checked bounds}\label{sec:protected-method}
The owner quantizes points and centers with input scale $S$, producing integers $\widetilde x$ and $\widetilde c_j$. The workload uses $S=4096/\max(\max|X|,\max|\mathcal C_0|,10^{-12})$, where $\mathcal C_0$ is the deployed set; later centers use the same scale and are clipped to the approved coordinate range. The basis scale is $T=4096$, and $B_U=\operatorname{round}(TU^\top)$. For $w=\widetilde x-\widetilde c_0$, Stage~1 proves
\begin{align}
v&=B_Uw,& H&=Tv,\nonumber\\
R&=T^2w-B_U^\top v,& E&=\sum_{\ell=1}^d R_\ell^2.\label{eq:protected-integer}
\end{align}
The decoded coordinates are $H/(T^2S)$ and residual energy $E/(T^4S^2)$. The circuit proves this integer computation exactly. Rounding $B_U$ generally destroys exact orthonormality, so Eq.~\eqref{eq:head-energy} is an approximate original-space statement for the implementation. We distinguish compression error from input quantization in the evaluation.

For concreteness, suppose $|\widetilde x_\ell|,|\widetilde c_{0\ell}|\le B$, and let $L_j=\sum_\ell |(B_U)_{j\ell}|$. Valid bounds are
\begin{align}
|H_j|&\le 2BTL_j,\nonumber\\
|R_\ell|&\le 2B\left(T^2+\sum_j|(B_U)_{j\ell}|L_j\right).\label{eq:row-bounds}
\end{align}
Let $H_{\max}$ be the maximum head bound and $R_{\max}$ the maximum residual bound; then $0\le E\le E_{\max}=dR_{\max}^2$. The implementation derives these bounds from the public integer model, range-checks signed inputs and heads and unsigned energies, and rejects model/aggregate bounds unsafe for the BN254 scalar field. No orthonormality assumption is needed for exact execution of Eq.~\eqref{eq:protected-integer}.

\subsection{Stage 1: original-data binding}\label{sec:query-relation}
The owner commits an ordered, power-of-two catalog using salted Poseidon2 row-Merkle trees. Each original leaf binds its index and full integer row. The derived tree binds the same index and the joint row $(H_i,E_i)$. A context binds the original root $C_X$, statistic roots, and a hash of the complete public model: dimensions, scales, anchor/deployed centers, integer basis and method. Unused optional roots are retained by the versioned context format.

A Stage~1 proof for index $i$ opens both the original row under $C_X$ and the derived row under $C_H$, proves Eq.~\eqref{eq:protected-integer}, and enforces the approved ranges. The verifier accepts preprocessing only after checking one proof for every catalog index exactly once, under one context and pinned verification key. This coverage check prevents a prover from supplying derived statistics for only a convenient subset. A later query cannot substitute another row or another basis without changing the authenticated statement.

\subsection{Learned head for changing centers}\label{sec:reuse}
To support new center sets, we retain the deployed-center span and append principal components of training residuals, orthogonalizing them against that span (\path{code/commitjl/enriched_jl.py}). The main workload uses $a=15$ head coordinates and one energy scalar. Neither basis fitting nor $k$-means training uses the held-out audit rows. The approved basis and original catalog remain unchanged across queries. These statistics are sufficient for exact real-arithmetic distances to in-span centers; for arbitrary out-of-span centers they support the approximation below.

For a new center $c'_j$, define its head and residual using this same basis and anchor. The exact real identity is
\begin{align}
\|x-c'_j\|^2={}&\|h(x)-h(c'_j)\|^2+e(x)+e(c'_j)\nonumber\\
&-2\langle r(x),r(c'_j)\rangle.\label{eq:new-center-exact}
\end{align}
The reusable estimate $\widehat d_j(x)$ drops the last term. It is exact if $r(c'_j)=0$. In general,
\begin{equation}\label{eq:reuse-error}
|\widehat d_j(x)-\|x-c'_j\|^2|\le2\sqrt{e(x)e(c'_j)}.
\end{equation}
This follows from Cauchy--Schwarz and does not assume a data distribution. If $a_i$ minimizes the estimated distances and $j$ is any competitor, a sufficient real-arithmetic stability condition is
\begin{equation}\label{eq:reuse-certificate}
\widehat d_j(x_i)-\widehat d_{a_i}(x_i)>
2\sqrt{e(x_i)}\left(\sqrt{e(c'_j)}+\sqrt{e(c'_{a_i})}\right).
\end{equation}
This real-arithmetic condition alone does not apply to the rounded integer basis. The plain reuse circuit reports only an aggregate diagnostic margin. The optional certified relation below adds a rounding-aware original-space guarantee.

\subsection{Stage 2: one circuit for successive public center sets}\label{sec:reuse-circuit}
The verifier quantizes every public new center with the deployed scale, checks its coordinate range, and computes $(H'_j,E'_j)$ by the same integer rule. These values are public inputs to \path{code/zk/circuit/bound_reuse.go}. Their consistency with the intended full centers is a verifier obligation; the circuit alone accepts any bounded public center statistics.

For public sorted distinct indices $I$, the witness supplies authenticated $(H_i,E_i)$, paths, a salt and private assignments. The circuit opens each row, checks its ranges, and compares
\begin{equation}
s_{ij}=2\langle H_i,H'_j\rangle-\|H'_j\|^2-E'_j.
\end{equation}
It selects the greatest score with lowest-index tie breaking, equivalently the smallest $D_{ij}=\|H_i-H'_j\|^2+E_i+E'_j$. It proves the public outputs
\begin{align}
\widehat\Psi&=\sum_{i\in I}D_{i,a_i},\nonumber\\
M&=\sum_{i\in I}(D_{i,(2)}-D_{i,(1)}),\label{eq:reuse-outputs}
\end{align}
where $(1),(2)$ denote the two smallest distances, counting ties. Cost and margin decode by $T^4S^2$. The energy $E_i$ cancels in comparisons but remains in the reported cost. A large $M$ can coexist with a few unstable records and does not certify Eq.~\eqref{eq:reuse-certificate}.

For bounded comparisons, the implementation shifts each score by
$\Delta=6aH_{\max}^2+2E_{\max}$ and bounds it below $2\Delta$. Distances are bounded by $4aH_{\max}^2+2E_{\max}$. Range checks and a field bound on the subset aggregate prevent modular wraparound. With the catalog model, subset size and number of new centers fixed, the constraint system and Groth16 setup are reused for all new center sets. Changing these circuit shapes requires another setup.

\subsection{Rounded-integer accuracy certificates}\label{sec:integer-certificates}
In this subsection, $x$ and $c_j$ denote the committed quantized integer vectors, and $B=B_U$. Define the public Gram defect
\begin{equation}
g=\max_p\sum_q |(BB^\top-T^2I)_{pq}|.
\end{equation}
For each record and center, use Eq.~\eqref{eq:protected-integer} with the same fixed anchor. Put $u_i=\lceil\sqrt{E_i}\rceil$, $u'_j=\lceil\sqrt{E'_j}\rceil$, and
\begin{align}
Q_{ij}&=T^4\|x_i-c_j\|^2,\nonumber\\
\delta_{ij}&=2T^2u_i u'_j+g\|H_i-H'_j\|^2.\label{eq:integer-delta}
\end{align}
\begin{proposition}[Error envelope for a rounded basis]\label{prop:integer-envelope}
For every integer basis $B$, record $x_i$ and center $c_j$,
\begin{equation}\label{eq:integer-envelope}
T^2|D_{ij}-Q_{ij}|\le\delta_{ij}.
\end{equation}
\end{proposition}
\begin{proof}
Set $v=x_i-c_j$ and $z=Bv$. Exact expansion gives
\begin{align*}
D_{ij}-Q_{ij}={}&2\langle R_i,R'_j\rangle\\
&+z^\top(BB^\top-T^2I)z.
\end{align*}
Cauchy--Schwarz bounds the first term by $2u_i u'_j$. For the symmetric matrix $G=BB^\top-T^2I$, the inequality $2|z_pz_q|\le z_p^2+z_q^2$ gives $|z^\top Gz|\le g\|z\|^2$. Multiply by $T^2$ and use $H_i-H'_j=Tz$.
\end{proof}
This deterministic bound uses no distributional hypothesis, center-span assumption or exact orthonormality. The Gram term cannot be dropped even when a center has zero residual energy: $T=8$, $B=[7]$, $x=1$, anchor/center zero gives $D=3361$, whereas $Q=4096$.

Write $d_{ij}=T^2D_{ij}$ and let $a_i$ be the compressed first-index minimizer. The certified circuit computes a private Boolean $b_i$ that is one exactly when
\begin{equation}\label{eq:integer-stability}
d_{ij}-d_{i,a_i}>\delta_{ij}+\delta_{i,a_i}
\quad\text{for every }j\ne a_i.
\end{equation}
It publishes $C_I=\sum_{i\in I}b_i$. Strict inequality conservatively leaves ties uncertified. The circuit checks that private row radii are nonnegative bounded integers with $u_i^2\ge E_i$; the verifier computes center radii from the intended public centers. Honest execution uses ceiling roots. Larger admissible row radii only weaken the guarantee, so $C_I$ is a lower bound on correct assignments, not a claim of maximum certifiable coverage.

The circuit also computes public interval endpoints
\begin{align}
L_I&=\sum_{i\in I}\max\{0,\min_j(d_{ij}-\delta_{ij})\},\nonumber\\
U_I&=\sum_{i\in I}\min_j(d_{ij}+\delta_{ij}).\label{eq:integer-interval}
\end{align}
Taking minima in Eq.~\eqref{eq:integer-envelope} proves
\begin{equation}\label{eq:integer-cost-contract}
L_I\le T^6\Phi_I(\mathcal C)\le U_I.
\end{equation}
The quoted approximate cost $A_I=T^2\widehat\Psi$ lies in the same interval. If $L_I>0$ and $U_I-L_I\le\epsilon L_I$, then
\begin{equation}
\frac{|A_I-T^6\Phi_I(\mathcal C)|}{T^6\Phi_I(\mathcal C)}\le\epsilon.
\end{equation}
Both endpoints decode by $T^6S^2$. They bound the quantized optimum, not a floating optimum before input rounding.

The implementation in \path{code/zk/circuit/bound_reuse_certificate.go} reuses the authenticated head/energy catalog and checks all error and interval arithmetic inside Stage~2. It derives conservative integer ranges from the public model and rejects unsafe comparison or aggregate bounds before compilation. It does not represent negative interval candidates as unconstrained field residues. Public center radii, certified count and interval endpoints extend the public statement; individual flags remain private. Neither a large aggregate margin nor high empirical agreement substitutes for this certificate.

An auditor specifies a minimum certified fraction $\alpha$ and maximum cost precision $\epsilon$, checks the public task and every required proof, and accepts a set only if $C_I\ge\alpha|I|$ and the interval meets the chosen cost condition. For disjoint batches, counts and endpoints add. The public-only policy checker uses integer comparisons and separates cryptographic validity from quality acceptance. Exact ties may prevent $C_I=|I|$ even when all computed assignments are correct. If the policy fails, the auditor refuses the approximate answer or requests public-center Full for that catalog; selective refinement is not implemented here.

\subsection{Exact public-center and fixed-task alternatives}\label{sec:baselines}
The public-center Full relation authenticates the original integer rows directly and takes every coordinate of every new center as a public input. It range-checks them, computes $\min_j\|x_i-c_j\|^2$ with first-index ties, and proves the exact sum. The witness contains the original rows, paths and private assignments. No statistic preprocessing or Stage~1 is needed. One compiled relation and Groth16 key serve every center set of the same catalog/model shape, subset size and $k$. The verifier builds the public witness from its independently approved full centers and ordered subsets. We implement this baseline in \path{code/zk/circuit/bound_full_reuse.go}, rather than repeatedly charging center-specific Full setup.

The cached baseline proves each original-space integer argmin and minimum cost in Stage~1, then authenticates and sums them in Stage~2. It stores two values per row and is the strongest tested control for one fixed set. Exact affine instead stores $\beta(x)=\|\widetilde x-\widetilde c_0\|^2$ and $k-1$ scores
\begin{equation}
f_j(x)=2\langle\widetilde x-\widetilde c_0,\widetilde c_j-\widetilde c_0\rangle-
\|\widetilde c_j-\widetilde c_0\|^2,
\end{equation}
with $f_0=0$. The nearest cost is $\beta-\max_j f_j$. Both reproduce integer Full exactly for their approved centers. Our center-specific cached and affine implementations repeat preprocessing for each new set. This is a property of the evaluated baselines: affine information could also be parameterized in a fixed spanning basis to support in-span changes.

PCA with energy stores 15 PCA coordinates plus a residual norm, applies the same new-center estimate, and uses the same plain or certified reuse circuit. Its basis uses only training rows; the proof establishes the approved transformation, not correctness of PCA training. The earlier historical Full circuit embeds center values and has center-specific setup. We distinguish those extrapolated historical totals from actual complete-catalog public-center Full measurements.
